\documentclass[11pt, a4paper]{article}

\usepackage[utf8]{inputenc}
\usepackage{amsmath, amssymb}
\usepackage{geometry}
\usepackage{hyperref}

\geometry{margin=1in}

\title{\textbf{Project Report: Iterative MCTS-SINDy Hybrid Architecture for Equation Discovery}}
\author{A Detailed Report Plan}
\date{\today}

\begin{document}

\maketitle

\begin{abstract}
    % Summary of the paper.
    % 1. Problem statement: Identifying governing equations from data.
    % 2. Baseline models: SINDy (requires predefined libraries) vs MCTS (combinatorial explosion, credit assignment).
    % 3. Our Contribution: An Iterative greedy MCTS-SINDy pipeline augmented by SHUSS (Sequential Halving).
    % 4. Results: 50x speedup on simple physics, success on Chaotic Systems (Lorenz), and exploration of algorithmic depth limitations on heavily nested compound functions.
\end{abstract}

\section{Introduction}
% - Context: Data-driven discovery in physics.
% - Limitations of pure Symbolic Regression (Genetic Algorithms/MCTS) $\to$ hard to scale, local minima.
% - Limitations of PySINDy (Sparse Identification of Non-linear Dynamics) $\to$ requires human-curated feature libraries, fails when basis functions like $\sin(x^2)$ are missing.
% - Proposed approach: Use MCTS to dynamically generate the SINDy library one feature at a time.
% - Outline of the report.

\section{Background and Foundations}
\subsection{Sparse Identification of Non-linear Dynamics (SINDy)}
% - How STLSQ regression works to enforce sparsity.
% - The matrix form: $\dot{X} = \Theta(X)\Xi$.

\subsection{Monte Carlo Tree Search (MCTS) for Syntax Trees}
% - State, Action space defined by Context-Free Grammars (CFG).
% - The UCT formula for balancing Exploration vs. Exploitation.

\subsection{Multi-Armed Bandits and Sequential Halving (SHUSS)}
% - Brief introduction to SHUSS (Sequential Halving Applied to Trees).
% - How allocating specific rollout budgets helps prune bad branches deterministically compared to flat MAB allocations.

\section{The Core Challenge: Generating the Library (Feature Entanglement)}
% - INITIAL ATTEMPT: Generating the ENTIRE SINDy dictionary $\Theta(X)$ via one massive MCTS string.
% - The Credit Assignment Problem: If the MCTS proposes $[\sin(y), \exp(x), \cos(x)]$ and only $\cos(x)$ is physically relevant, SINDy still yields a low MSE. The MCTS backpropagates a high reward to $\sin(y)$ and $\exp(x)$, hopelessly entangling good and bad features.
% - Result: Total failure to converge on true dynamics cleanly without fitting to noise.

\section{Proposed Architecture: Iterative MCTS}
% - THE SOLUTION: Shift from a Global library generation to a greedy, Iterative search (similar to Orthogonal Matching Pursuit).
\subsection{The Greedy Additive Loop}
% - Formally define the loop: MCTS only generates ONE single mathematical expression $\phi$.
% - Once found, $\phi$ is permanently locked into a global dictionary $\Theta_{\text{locked}}$, and the model restarts the search focusing purely on the residual error.

\subsection{State Representation and Grammar}
% - Define the production rules ($S \to \dots$).
% - Show how the grammar restricts the search space to a single feature and enforces valid mathematical topology.

\section{Solving the MCTS Bottlenecks}
% Detail the major engineering/mathematical changes implemented to stabilize the model.

\subsection{Reward Design: BIC and Occam's Razor}
% - Moving away from unstable, raw Mean Squared Error (MSE).
% - Implementing the Bayesian Information Criterion (BIC): $BIC = n \ln(MSE) + k \ln(n)$.
% - The Hard Threshold: Terminating the iterative loop if the new feature fails to decrease the BIC by a significant margin (protecting against micro-noise overfitting).

\subsection{Reward Stability: Average vs. Max Tracking}
% - Issue: Outlier rollouts generating lucky overfit sequences caused "Max Reward" nodes to permanently lock in bad behavior.
% - Solution: Transitioning to Incremental Average Return bounds to stabilize the UCT scores.

\subsection{Simulation Tuning: Eliminating the Monte Carlo Blind-Walk}
% - Issue: Uniform Random Simulation (rollouts) generated 99% syntactic garbage.
% - Solution: Replacing flat rollouts with the SHUSS Bandit Tournament, evaluating partial subsets, and continuously pruning the worst 50% of the tree during simulation.
% - Result: 50x speedup in convergence (100 episodes vs 5000 episodes).

\section{Experimental Results \& Benchmarking}
% Validating the model across a spectrum of difficulty.

\subsection{Linear Baseline: Damped Harmonic Oscillator}
% - $dx/dt = -0.1x + y$.
% - Showcase perfect accuracy and extreme convergence speed due to SHUSS.

\subsection{Chaotic Systems: The Lorenz Attractor}
% - $dz/dt = x \cdot y - 2.66 \cdot z$.
% - Highlighting the model's capacity to resolve coupled non-linear dimensions ($x \cdot y$) while inherently rejecting massive polynomial libraries that typically overfit chaotic spaces.

\subsection{The Curse of Dimensionality: Deep Nesting Bounds}
% - The boundary benchmark: $\dot{y} = x + \sin(x^2 \cdot z \cdot y) + \cos(\sin(z \cdot y))$.
% - Discussing the 128-configuration Grid Search.
% - Finding: SHUSS accurately prunes shallow functions, but violently fails on depths $d > 8$ because partial combinations like $\sin(x...)$ yield *worse* intermediate residuals than simple errors $x$.
% - MCTS Success relies strictly on Monte Carlo Lottery/Luck at extreme composition depths.

\section{Conclusion and Future Work}
% - Summary: We successfully integrated MCTS and SINDy by mapping the combinatorial problem into a greedy additive feature search, and resolved scaling issues with SHUSS.
% - Future limits: Unguided tree search hard-caps around derivation depths of 8 without exponential compute scaling.
% - Future Extensions: Replacing Uniform Random Rollouts in SHUSS with a Graph Neural Network (GNN) Policy Prior to explicitly guide physics-aware syntax generation.

\end{document}
